\begin{abstract}
Vision-Language-Action (VLA) models have demonstrated strong multi-modal reasoning capabilities, enabling direct action generation from visual perception and language instructions in an end-to-end manner. However, their substantial computational cost poses a challenge for real-time robotic control, where rapid decision-making is essential. This paper introduces VLA-Cache, a training-free inference acceleration method that reduces computational overhead by adaptively caching and reusing static visual tokens across frames. Exploiting the temporal continuity in robotic manipulation, VLA-Cache identifies minimally changed tokens between adjacent frames and reuses their cached key-value representations, thereby circumventing redundant computations. Additionally, to maintain action precision, VLA-Cache selectively re-computes task-relevant tokens that are environmentally sensitive, ensuring the fidelity of critical visual information. To further optimize efficiency, we introduce a layer adaptive token reusing strategy that dynamically adjusts the reuse ratio based on attention concentration across decoder layers, prioritizing critical tokens for recomputation. Extensive experiments on two simulation platforms (LIBERO and SIMPLER) and a real-world robotic system demonstrate that VLA-Cache achieves up to \textbf{1.7× speedup} in CUDA latency and a \textbf{15\% increase} in control frequency, with negligible loss on task success rate. The code and videos can be found at our project page: \url{https://vla-cache.github.io}.
\end{abstract}

